% Capitulo: Antes e depois da IA no meu trabalho
\chapter{Antes e depois da IA no meu trabalho}
\label{cap:antes-depois}

\rascunho{
  \textbf{Ideia do capítulo.} Comparação direta do mesmo fluxo de trabalho nos dois momentos.
  \begin{itemize}
    \item Estimativa e refinamento: o que mudou no tamanho das tarefas.
    \item Escrita de código e revisão: mais rápido para escrever, mais importante revisar.
    \item Testes: antes, TDD ficava para quando sobrasse tempo e as regras eram conferidas
      com o stakeholder a cada entrega; depois, os testes vêm antes do modelo e a
      validação com o stakeholder acontece uma vez, sobre os casos de teste (ver \veja{sec:tdd}).
    \item O diferencial migrou do técnico para o negócio (ver o capítulo~\ref{cap:negocio}).
    \item Custo: da licença por usuário ao consumo em créditos e tokens, e o gerente
      acompanhando consumo de forma contínua (ver o capítulo~\ref{cap:custo-consumo}).
    \item Documentação e transferência de conhecimento.
    \item O que piorou ou exigiu novos cuidados.
    \item Usar a caixa \codigo{\textbackslash antesdepois} ao longo do capítulo.
  \end{itemize}
}

\section{Estimativa e refinamento}
\label{sec:ad-estimativa}

\pendente{escrever}

\section{Código e revisão}
\label{sec:ad-codigo}

\pendente{escrever}

\section{Documentação}
\label{sec:ad-documentacao}

\pendente{escrever}

\section{O que exigiu novos cuidados}
\label{sec:ad-cuidados}

\pendente{escrever}
